\chapter{Evaluación de la Calidad Industrial}

\section{Real Decreto 2200/1995: Acreditación, certificación y laboratorios de ensayo}
\subsection{Acreditación}
La acreditación es la herramienta establecida a escala internacional para \textbf{generar confianza} sobre la actuación de un tipo de organizaciones muy determinado que se denominan de manera general \textbf{Organismos de Evaluación de la Conformidad} y que abarca a los \textit{Laboratorios de ensayo, Laboratorios de Calibración, Entidades de Inspección, Entidades de certificación y Verificadores Ambientales}

\subsubsection{¿Por qué existe?}
Para lograr la \textbf{CONFIANZA Y CREDIBILIDAD} que el mercado y la sociedad en general tiene de los \textbf{evaluadores que realizan las evaluaciones}, por lo que es preciso establecer un mecanismo independiente, riguroso y global que \textbf{garantice la competencia técnica} de dichos evaluadores y su sujeción a normas de carácter internacional.

\subsubsection{Entidades de Acreditación}
\textit{``Se reconoce y designa a la Entidad Nacional de Acreditación (ENAC) como Entidad de Acreditación de las establecidas en el Capítulo 2''.} - RD 2200:1995 (modificado por el RD 338:2010)

\paragraph{CAPÍTULO 2. ENTIDADES DE ACREDITACIÓN.} \textit{``Son entidades privadas sin ánimo de lucro, que se constituyen con la finalidad de acreditar en el ámbito estatal, a las entidades de certificación, laboratorios de ensayo y calibración y entidades auditoras y de inspección que actúen en el ámbito voluntario de la CALIDAD, así como a los Organismos de Control que actúen en el ámbito reglamentario y a los Verificadores medioambientales, mediante la verificación del cumplimiento de las condiciones y requisitos técnicos exigidos para su funcionamiento.''}

\subsubsection{¿Por qué existen los Evaluadores de la Conformidad?}
El objetivo principal de la actuación de los organismos de evaluación de la conformidad es el de demostrar a la sociedad (Autoridades, empresas y consumidores en general) que los PRODUCTOS Y SEVRICIOS puestos a su disposición son conformes con ciertos requisitos relacionados generalmente con su CALIDAD y la SEGURIDAD.

Dichos requisitos pueden estar establecidos por la Ley y tneer por tanto CARÁCTER REGLAMENTARIO o estar especificados en Normas, especificaciones u otros documentos de CARÁCTER VOLUNTARIO.

\subsubsection{¿Cuál es el proceso de una acreditación?}


\subsection{Certificación}
La Acreditación reconocía la competencia técnica de una organización para la realización de ciertas actividades, con las siguientes ventajas:
\begin{itemize}
    \item Declara que los organismos acreditados son competentes e imparciales;
    \item Les permite conseguir la aceptación de sus prestaciones a nivel internacional y el reconomiento de sus competencias.
    \item Evita a las empresas exportadoras los reiterados controles que deben pasar para tener acceso a los mercados internacionales
    \item Acceso a los mercados internacionales
    \item Establece y promueve la confianza a nivel nacional e internacional al comprobar la competencia de los operadores
\end{itemize}

La certificación es la acción llevada a cabo por una entidad reconocida como independiente de las partes interesadas, medianta la que se dispone de la confianza adecuada en que un producto, proceso o servicio debidamente identificado, es conforme con una norma u otro documento normativo especificado.

\subsubsection{¿Por qué se hace?}

El concepto de calidad ha evolucionado fuertemente a lo largo del siglo XX. se ha avanzado hacia la satisfacción de los requerimientos de los clientes, siguiendo luego por la adaptación para el costo, que implica incluir el aspecto económico.

La evolución de calidad tiene cuatro etapas:

\subsubsection{¿Qué se certifica?}
Actualmente podemos distinguir dos grandes grupos:

El HACCP es una herramienta para evaluar peligros y establecer sistemas de control centrados en la prevención. Puede aplicarse en toda la cadena alimentaria, desde el productor primario hasta el consumidor final.

\subsubsection{¿Cómo se hace?}

\subsubsubsection{Certificación de Sistemas de Calidad}
En esta categoría podemos encontrar la certificación de sistemas en base a las normas ISO 9000, ISO 14000, OSHAS 18000 u otras:
\begin{itemize}
    \item La ISO 9001:2000 y la ISO 14001:2004 dan los requisitos genéricos para los sistemas de gestión, no los requisitos para los productos específicos o los servicios.
    \item Las certificaciones de la ISO 9001:2000 y de la ISO 14001:2004 no son certificaciones del producto o garantías del producto.
    \item Se deben tomar las precauciones pertinentes de cualquier referencia a ella en la información relacionada con el producto, incluyendo los anuncios, o en cualquier otro medio, para evitar de dar la impresión que son certificaciones del producto o garantías del producto.
    \item En detalle, las marcas de certificación de la ISO 9001:2000 y de la ISO 14001:2004 NO deben ser exhibidas en rpoductos, en etiquetas del producto, en el producto que empaqueta, o de ninguna manera que se pueda interpretar como denotar conformidad del producto.
\end{itemize}

\subsubsubsection{Certificación de Productos}
\paragraph{Modelos de Certificación de Producto ISO / CASCO}

Se describen los ocho modelos de Certificación existentes y la diferencia entre ellos, partiendo por el Ensayo de Tipo (Modelo Nº 1), pasando por la Inspección de Lote (Modelo Nº 7), para terminar en el Modelo Nº 8, que corresponde a la Inspección 100\%.




\subsection{Laboratorios de Ensayo}


\section{Directivas de Nuevo Enfoque. Evaluación de la conformidad}

\section{El marcado CE}

\section{Campañas de control reglamentario}